\section{An integral involution projector and the remaining normal-length gate}
\label{sec:kt-projector}

This section separates a new integral conormal criterion from a
source-located limitation of the existing point-counting argument.
The subsequent KT-ODD-CONORMAL-SELMER-EQUIVALENCE identifies
that conormal with the dual of the original strict Selmer
module.  Its minor criterion is therefore an equivalent
formulation, with no arithmetic gate reduction claimed.
All source references use the pinned commit of KT-ONECHAR-CLOSED.
These are mathematical proofs and source inspections; no Lean
compilation is claimed.  The hypotheses and notation of
KT-ADJOINT-ACTUAL-HOST-POINT are retained.  In particular $O$ is
any finite coefficient DVR over $\mathbf Z_\ell$, $\ell>3$,
$e\geq2$ is even, $\eta=\chi_{-3}$, $\Psi=\varepsilon_\ell^e$,
and $(\ell-1)\nmid2e$.

\subsection{Part I: what the present length proof actually detects}

Write $\mathcal A_j$ for the source's torus weight spaces and
$x:A\to O$ for its base augmentation.  The literal definitions in
\texttt{CommutativeAlgebra/ResMapBijective.lean:245--268} are
\[
 \begin{split}
 \operatorname{satSq}(x)
   &=\{f:\text{ some }0\ne c\in O\text{ has }cf\in(\ker x)^2\},\\
 L_{-1}&=\operatorname{im}\bigl(\mathcal A_{-1}\cap\ker x
                     \longrightarrow A/\operatorname{satSq}(x)\bigr),\\
 C_A&=L_{-1}/\langle F_i\bmod\operatorname{satSq}(x)\rangle_O.
 \end{split}                                                    \tag{KP.1}
\]
Here $F_i$ are the negative weight-one equations.  The quantity
\texttt{cotLen} is $\operatorname{length}_O C_A$.  Saturation
makes the ambient lattice $L_{-1}$ torsion-free; it does not make
its quotient $C_A$ torsion-free.  The latter integral torsion is
precisely the quantity bounded.

The following is the exact chain used in the archived proof.
In the middle four entries, filenames are relative to
\texttt{CommutativeAlgebra/GradedModulesLength/}.
\begin{center}
\begin{tabular}{p{0.34\textwidth}p{0.59\textwidth}}
\textbf{Source locator} & \textbf{Operation relevant to the projector}\\\hline
\texttt{UnitaryThreeGroups/ModelLower.lean:296--369,374--398}
 & Constructs the normalized dual-number point from the strict
 Selmer cocycle, then lifts its derivation to the ambient model.\\
\texttt{Mprime.lean:38--80}
 & Injects the torsion input into
 $\operatorname{Hom}_O(C_A,O/(a))$.  It does not inject $C_A$
 into a generic point algebra.\\
\texttt{MainCotSpan.lean:32--70}
 & Takes the Smith normal form of the integral linearized
 negative equations, cuts the zero-weight variables, and invokes
 the Rees/Euler-characteristic estimate.\\
\texttt{ToDegInjective.lean:187--219}
 & Introduces $F_i^2$ and $E_i^2$ in the Rees sections and the
 auxiliary ideal $I_2$ in (KP.2).\\
\texttt{TwoSigned.lean:17--95}
 & Bounds localized lengths of the auxiliary quotient by
 $2^{r+1+r'}$ times generic orbit valuations.\\
\texttt{AlgMemCoef.lean:92--143}
 & Compares the lower contribution
 $2^{r+r'}\sum_i c_i$ with that upper bound, yielding the factor
 $2$ in the final estimate.\\
\texttt{EtHolds.lean:71--103,123--137}
 & Detects elements by characteristic-zero points only after
 every nonzero element of $O$ is assumed invertible.\\
\end{tabular}
\end{center}
The equations $F_i,G_i,E_i,E'_i,Z_i$ have weights
$-1,-2,1,2,0$, respectively.  The exact auxiliary ideal is
\[
        I_2=(F_i^2,\ G_i,\ E_i^2,\ E'_i,\ Z_i).              \tag{KP.2}
\]
Thus even if a chosen $F_i$ is odd under an involution,
its square is even, and the class of $F_i$ can remain an odd
nilpotent in the auxiliary quotient.  The localized-length
bound explicitly allows this thickening.  The generic reduced
orbit bound by one occurs before multiplying by the powers of
two which account for (KP.2).

More precisely, \texttt{EtHolds.et\_holds} assumes
\texttt{\(\forall c\ne0,\) IsUnit (algebraMap O C c)}.
Its proof constructs an algebra structure over $\operatorname{Frac}O$
at lines 84--93 and obtains finiteness and reducedness over that
field at lines 99--103.  This is not a flatness or reducedness
assertion for the integral host ring.  None of the displayed
length statements has an involution as an argument.  The audited
chain therefore supplies no injection of the integral odd
conormal into a torsion-free reduced trace algebra.  This is a
limitation of the proposed refinement, not a defect in the
stated source inequality.

\begin{lemma}[KT-FIXED-GENERIC-POINTS]
Let $B$ be a finite torsion-free $O$-algebra with reduced generic
fiber, and let $\kappa$ be an $O$-algebra involution.  If every
geometric generic point is fixed by $\kappa$, then $\kappa=1$ on
$B$.  The analogous conclusion for the integral odd cotangent
of an arbitrary augmented complete local $O$-algebra does not
follow from its reduced generic point set.
\end{lemma}
\begin{proof}
The finite reduced algebra $B\otimes_O\operatorname{Frac}O$
is a product of finite separable fields, since the coefficient
field has characteristic zero.  Its geometric points separate
elements.  Hence $b-\kappa(b)$ vanishes in the generic fiber
for every $b\in B$, and torsion-freeness gives $b=\kappa(b)$.

For the second assertion let $\varpi$ be a uniformizer and, for
any $n\geq1$, set
\[
 R_n=O[[u]]/(\varpi^n u,u^2),\qquad
 x(u)=0,\qquad \kappa(u)=-u.                                \tag{KP.3}
\]
Give $u$ weight $-1$.  As an $O$-module,
$R_n=O\oplus (O/\varpi^n)u$.  Its generic fiber and its maximal
torsion-free reduced quotient are both $\operatorname{Frac}O$
and $O$, respectively; the unique generic point is fixed.
Nevertheless
\[
       (\ker x/(\ker x)^2)_{-1}^{-}
                    =(O/\varpi^n)u                           \tag{KP.4}
\]
has length $n$.  The example is a counterexample to the
generic-reduction inference alone; it is not claimed to satisfy
all the source hypotheses \texttt{M1--M4}.
\end{proof}

Even in a flat reduced algebra, an individual fixed point
can have nonzero integral odd cotangent when other generic
points are present.  For example
\[
 R'_n=O[[u]]/(u(u^2-\varpi^{2n})),\qquad \kappa(u)=-u          \tag{KP.5}
\]
is finite free of rank three over $O$, and reduced.  The generic
points $u=\varpi^n$ and $u=-\varpi^n$ are interchanged, while
$u=0$ is fixed.  At the augmentation $u=0$ its odd cotangent is
$O/(\varpi^{2n})$.  Indeed the only linear relation is
$-\varpi^{2n}u$, so its length is $2n$.  Reducedness follows
from the embedding into its reduced generic fiber, using
torsion-freeness.  This second example distinguishes a fixed
point from a branch with trivial integral normal conormal.
It does not refute a possible length bound counting only
nonfixed generic points: the two interchanged points could
account for the length $2n$ in this example.

\subsection{Part II: an involution on the actual normalized ring}

Keep $J$ and $K$ from (HS.1) and (HS.4), and put
$D=\operatorname{diag}(-1,1,-1)$.  The matrix automorphism
\[
 \kappa(M)=D\theta(M)D=K M^{-t}K                             \tag{KP.6}
\]
extends to $GG_3$ by leaving its sign component unchanged.
It commutes with $\theta$, fixes $r_0$, and commutes with
conjugation by $T(t)=\operatorname{diag}(t,1,t^{-1})$.
We first verify the local condition needed to use (KP.6) on
the source's integral functor.

\begin{lemma}[KT-ORDINARY-INTEGRAL-COMPONENT]
Under the hypotheses above and the carried local component
theorem, the source's local quotient $U.R/I_t$ represents the
framed ordinary functor with prescribed ordered inertial
characters $(\varepsilon_\ell^e,1,\varepsilon_\ell^{-e})$.
It is a power series ring in twelve variables over $O$.
Consequently $I_t$ is stable under $\kappa$.
\end{lemma}
\begin{proof}
The proof of KT-ORDINARY-ON-BASE-COMPONENT already establishes
formal smoothness of the fixed-inertia Borel lifting functor,
including unramified nonsplit $\eta$.  Its relative tangent
dimension is nine, as follows.  The twisted adjoint Borel Lie
algebra $\mathfrak b$ has dimension six and $H^2(G,\mathfrak b)=0$.
Local Euler characteristic over $G=G_{\mathbf Q_\ell}$ gives
$h^1(\mathfrak b)-h^0(\mathfrak b)=6$.  Therefore the framed
cocycle space $Z^1(G,\mathfrak b)$ has dimension twelve.
The map to $H^1(I,\mathfrak t)^{G/I}$ is onto by the proof of
(HS.8).  Since $\dim\mathfrak t=3$, $H^2(G,\mathfrak t)=0$,
and $\dim H^1(G/I,\mathfrak t)=h^0(\mathfrak t)$, the latter
space has dimension three.  Prescribing inertia leaves nine
parameters.  In the nonsplit case $h^0(\mathfrak t)=1$, whereas
in the split case it is three; the cancellation just used works
in both cases.  Thus $R_B\simeq O[[z_1,\ldots,z_9]]$.

Choose an inertia element $\tau$ for which the three residual
eigenvalues $\bar\Psi(\tau),1,\bar\Psi(\tau)^{-1}$ are distinct;
the stated GD condition ensures such an element.  The
characteristic polynomial of the unrestricted universal
matrix $r(\tau)$ has three simple residual roots.  Hensel
factorization in the unrestricted complete local ring itself
therefore gives three ordered roots and their spectral
idempotents.  The resulting ordered eigenflag lies in the big
cell around the standard flag.  For an ordinary lift it is the
unique invariant flag in that neighborhood: an invariant flag
is stable under $r(\tau)$ and its residual ordering specifies
which spectral lines it contains.

The big cell is uniquely parametrized by a lower unipotent
matrix congruent to one, with three coordinates.  Consequently
every framed ordinary lift is uniquely a conjugate of a Borel
lift by such a matrix.  This gives a representing ring
\[
              R_{\rm ord}=R_B[[z_{10},z_{11},z_{12}]].        \tag{KP.7}
\]
It also shows that this is a closed subfunctor of the
unrestricted framed functor.  Indeed use its canonical Hensel
eigenflag, require every $r(g)$ to preserve it, and require the
three diagonal inertia characters to have the prescribed
values.  These are closed matrix equations in the universal
ring.  For the sign $-1$ action, preservation means preservation
under the semidirect action on flags; conjugation into
$B\rtimes\{\pm1\}$ is the same condition because $\theta$
preserves the upper Borel.  Thus $U.R\to R_{\rm ord}$ is
surjective, with prime kernel $P$.

KT-ORDINARY-ON-BASE-COMPONENT gives $I_t\subset P$.  The carried
theorem \texttt{component\_basic} states that $I_t$ is prime and
$\dim U.R/I_t=13$; see
\texttt{OrdinaryComponentUniquePrimeDimThirteenNearBase.lean:11--35}.
Equation (KP.7) gives $\dim U.R/P=13$.  A strict inclusion of
these primes would append one more prime to a chain of length
thirteen in $U.R/P$, a contradiction.  Hence $P=I_t$.

Finally $\kappa$ preserves $B$, reverses and inverts its three
diagonal characters, and therefore preserves their prescribed
ordered list.  It carries a congruent conjugating matrix to
another such matrix.  It thus preserves the represented
ordinary functor, proving the asserted stability of $I_t$.
\end{proof}

\begin{lemma}[KT-EQUIVARIANT-MARKED-SLICE]
The actual source ring
\[
 S=\operatorname{Univ}\bigl(\operatorname{NormCond}(c,g_1)
                        (\operatorname{BoxCond}(\mathrm{Kill},I_t))\bigr)
\]
admits an $O$-algebra involution $\kappa_S$, commuting with the
source torus action and fixing its base augmentation.  On
unmarked representations it is (KP.6); on the marked slice it
is this operation followed by a functorial normalization.
The ordinary adjoint points constructed above admit
$\kappa_S$-fixed representatives.  On the rational strict
Selmer dual-number directions, $\kappa_S$ acts by $-1$.
\end{lemma}
\begin{proof}
The local lemma, $\kappa(r_0)=r_0$, and $\kappa(1)=1$ show that
the unnormalized box conditions are stable.  Congruent
conjugation preserves them by the source's local conjugacy
theorem and the definition of the killed subgroup.

It remains to treat the actual marking, rather than assume
that its linearization suffices.  Work first with
$r(c).\mathrm{mat}=1$, as in the source's normalized functor,
and set $M=r(g_1).\mathrm{mat}$.  Since $\eta(g_1)=1$, the two
remaining marking equations are
\[
 Z_1(M)=M_{12}-M_{23}=0,\qquad
 Z_2(M)=M_{21}-M_{32}=0.                                    \tag{KP.8}
\]
This is the literal marking after
\texttt{MatGstar} and \texttt{NormEx.anti\_shape}; indices here
are one-based.  Define genuinely odd functions
\[
           F_i(M)=\frac{Z_i(M)-Z_i(\kappa M)}2.              \tag{KP.9}
\]
They have torus weights $1$ and $-1$ and satisfy
$F_i(\kappa M)=-F_i(M)$.

Put $X=E_{12}-E_{23}$, $Y=E_{21}-E_{32}$, and use the Cayley
cross-section
\[
 W=xX+yY,\qquad
 H(x,y)=(1-W/2)^{-1}(1+W/2),\qquad x,y\in\mathfrak m_R.
                                                                  \tag{KP.10}
\]
Both $X$ and $Y$ lie in $\mathfrak{so}(J)$.
Consequently $H J H^t=J$, so conjugation by $H$ preserves
$r(c).\mathrm{mat}=1$.  The cross-section has the identities
\[
 \begin{split}
 H(x,y)^{-1}&=H(-x,-y),\\
 \kappa H(x,y)&=H(-x,-y),\\
 T(t)H(x,y)T(t)^{-1}&=H(tx,t^{-1}y).
 \end{split}                                                    \tag{KP.11}
\]
The first identity, in particular, stays within the same
cross-section; an arbitrary ordered product of root groups
would not have this property.

At $M_0=\operatorname{diag}(q,1,q^{-1})$, the Jacobian in
$(x,y)$ of either pair
$Z_i(HMH^{-1})$ or $F_i(HMH^{-1})$ is
\[
               \begin{pmatrix}q^{-1}-q&0\\0&q-q^{-1}\end{pmatrix}.
                                                                  \tag{KP.12}
\]
For $Z$ this follows by differentiating $[xX+yY,M_0]$.
For $F$ the same derivative follows because these gauge
directions are $\kappa$-odd and $M_0$ is fixed.  The residual
value $q=\bar\Psi(g_1)$ satisfies $q^2-1\in k(O)^\times$
by \texttt{D.Marked(g1)}.  Thus the formal implicit function
theorem gives a unique $(x,y)\in\mathfrak m_R^2$ making either
pair of marking functions zero.  Equivalently one may solve
successively modulo every power of $\mathfrak m_R$ using the
unit Jacobian, and take the unique limit.  This also proves
functoriality in every complete coefficient algebra.

Let $N_F$ send the old $Z$-slice to the new $F$-slice using
this gauge, and let $N_Z$ send the new slice to the old.
They are mutual inverses: if $r'=H\cdot r$, then
$H^{-1}=H(-x,-y)$ is a gauge in the same cross-section taking
$r'$ back to its required old marking; uniqueness forces it
to be the inverse normalization.  The same argument applies
in the other direction.  Formula (KP.11) and uniqueness show
that $N_F$ is torus-equivariant, and that the $F$-normalization
on the $c$-normalized functor is $\kappa$-equivariant.
Transport the genuine involution of the $F$-slice through
$N_Z\circ\kappa\circ N_F$.  This proves the functorial
involution on the old slice, and representability gives
$\kappa_S$.  The base lift needs no gauge, so its augmentation
is fixed.

The original global adjoint lift in (HS.14) is literally fixed by
(KP.6): its matrices are $A(g)D^{j(g)}$, where $A(g)$
preserves $K$, and both $A(g)$ and $D$ are fixed by $\kappa$.
Its matrix at $c$ is already one.  The local conjugation in
(HS.15) was used to verify its ordinary local property;
we do not replace the global representative by that locally
conjugated one.  The original lift is therefore on the
$F$-slice.  Applying $N_Z$ gives a $\kappa_S$-fixed point on
the actual old slice, with all the host conditions unchanged.

Finally the source's rational strict Selmer cocycles have
the literal shape $\gamma(g)=t(g)(E_{21}-E_{32})$ and satisfy
$t(c)=t(g_1)=0$; see
\texttt{ShaInjGen.lean:101--129,153--187,216--288}.
The differential of $\kappa$ is $X\mapsto-KX^tK$, which sends
$E_{21}-E_{32}$ to its negative.  The corresponding
dual-number lift has $M=M_0$ exactly, so both marking gauges
are one.  The transported involution therefore sends its
derivation to its negative, as claimed.
\end{proof}

\subsection{The exact unsquared conormal inequality}

Let $x_0:S\to O$ denote the base augmentation, put
$J_0=\ker x_0$, and write
\[
       C^-_{-1}=\bigl(J_0/J_0^2\bigr)_{\text{weight }-1}^{\kappa_S=-1}.
                                                                  \tag{KP.13}
\]
Unlike (KP.1), this is the ordinary integral cotangent of the
actual ring.  No saturation or reduced quotient is taken.
It is a finite $O$-module, although finite \emph{length} is an
additional assertion when needed below.  The finite module
decomposes by torus weights and by the two idempotents
$(1\pm\kappa_S)/2$.

\begin{theorem}[KT-ODD-CONORMAL-GATE]
For every $0\ne b\in O$, the rational strict Selmer torsion
chunk $\mathcal X[b]$ used in \texttt{sha\_inj} injects into
\[
               \operatorname{Hom}_O(C^-_{-1},O/(b)).         \tag{KP.14}
\]
Here $\mathcal X$ is exactly the rational strict Selmer module
for the character $(\Psi)^{-1}$ in the source's convention.
In particular, if $C^-_{-1}$ has finite length, then
\[
           \operatorname{length}_O\mathcal X
                   \leq\operatorname{length}_O C^-_{-1}.    \tag{KP.15}
\]
Choose a finite torus- and involution-equivariant presentation
of $S$ over $O$.  Let $u_1,\ldots,u_n$ be its odd variables of
weight $-1$, and let $f_1,\ldots,f_m$ be its odd defining
equations of weight $-1$.  Write
\[
 f_j=\sum_{i=1}^n b_{ij}u_i+\text{terms of degree at least two},
       \qquad B_-=(b_{ij})\in M_{n\times m}(O).              \tag{KP.16}
\]
Then $C^-_{-1}=\operatorname{coker} B_-$.  If $\Delta$ is any
nonzero $n$-rowed minor, then
\[
       \operatorname{length}_O\mathcal X
                  \leq\operatorname{length}_O C^-_{-1}
                  \leq v_O(\Delta).                         \tag{KP.17}
\]
A unit maximal minor therefore proves vanishing.  If the
matrix has full rank, the best bound in (KP.17) is exactly
the valuation of its maximal-minor ideal.  For $n=0$ the
empty minor is one and the module is zero.
\end{theorem}
\begin{proof}
The explicit proof of \texttt{ShaInjGen} gives an injective
map from $\mathcal X[b]$ to the normalized cocycles just used.
Although the abstract \texttt{sha\_inj} statement permits an
annihilator $a$, its proof takes $a=b$; see
\texttt{ShaInj.lean:29--38}.  The construction
\texttt{ModelLower.main\_der} gives, before passing to the
ambient model, a derivation
$D:S\to O/(b)$ at $x_0$ with
\[
 D\bigl(r(g)_{ij}\bigr)=\Psi(g)^{\mathrm{wt}_1(j)}\gamma(g)_{ij}.
                                                                  \tag{KP.18}
\]
The unit factor in (KP.18) and the injectivity of the cocycle
map make the map to derivations injective.  Equivalently, two
derivations agreeing on all universal matrix entries agree
by uniqueness of the representing dual-number lift.  Each
derivation factors through $J_0/J_0^2$.  The cocycle has only
weight $-1$, so torus-equivariance, or (KP.18) applied to
torus conjugates, makes it zero on every other weight.
KT-EQUIVARIANT-MARKED-SLICE gives $D\circ\kappa_S=-D$.
Since two is a unit, it is zero on the even summand.  This
proves (KP.14).  The construction is additive; no unproved
compatibility with the saturated ambient quotient is used.

For a finite-length $O$-module $C$, the elementary divisor
decomposition shows
\[
\operatorname{length}_O\operatorname{Hom}_O(C,O/(b))
\leq\operatorname{length}_O C.
\]
The finite residue field
converts the injection into the same length inequality for
each torsion chunk, even if only its additive injectivity is
used.  The source's $\mathcal X$ is torsion, as in
\texttt{VlenLe.lean:13--42}; taking the increasing union of
its bounded torsion chunks proves (KP.15).

For completeness, an equivariant finite presentation is
obtained by choosing weight and involution eigenvectors in
$J_0/(J_0^2+\mathfrak m_OJ_0)$, lifting them to $J_0$, and
using the corresponding power series variables.  The torus
is diagonalizable and two is invertible, so the lifts can be
projected to the specified weight and sign.  The resulting
map from a finite power series ring is onto by complete
Nakayama.  Its ideal is closed and finitely generated;
apply the same weight/sign projections and Nakayama to
choose finitely many homogeneous eigen-equations.  The
formal torus action here is the one induced by the entry
weights; the Cayley change of slice is equivariant for it.

Modulo the square of the augmentation ideal, an equation
can affect the odd weight $-1$ variables only if it has that
same weight and sign.  Thus (KP.16) is exactly a presentation
of (KP.13).  Selecting a nonzero maximal minor gives a
square full-rank relation submatrix; its cokernel surjects
onto $\operatorname{coker}B_-$.  Smith normal form proves
the bound by $v_O(\Delta)$ and the equality for the
maximal-minor ideal.
\end{proof}

This is an exact integral conormal criterion, but no
matrix $B_-$ or unit minor has been evaluated in this run.
KT-ODD-CONORMAL-SELMER-EQUIVALENCE further proves that
it is equivalent to the original strict Selmer vanishing
problem.  In particular (KP.17) deletes no automorphic
branch from the existing trace-product bound and is not
counted as a stricter arithmetic gate.

\subsection{Unsquared normal equations along the adjoint branch}

The following exact algebra describes what a proposed
replacement by an arithmetic Jacobian would have to measure.
For any augmented ring with an involution as above, let
\[
 Q=(s-\kappa_S(s):s\in S),\quad B=S/Q,\quad N=Q/Q^2.         \tag{KP.19}
\]
Thus $B$ is the fixed-locus quotient, and $N$ its integral
normal conormal.  The involution is the identity on $B$
and is minus the identity on $N$.

\begin{lemma}[KT-UNSQUARED-NORMAL-JACOBIAN]
There is a canonical isomorphism
\[
       (J_0/J_0^2)^-\simeq N\otimes_{B,x_0}O.               \tag{KP.20}
\]
For an equivariant presentation
$S=O[[v_1,\ldots,v_a,u_1,\ldots,u_b]]/(f^+_j,f^-_j)$,
where $v$ is even and $u$ is odd, the $B$-module $N$ has the
presentation
\[
 B^{\{f^-_j\}}\xrightarrow{\,\mathcal J^-\,}B^b
              \longrightarrow N\longrightarrow0,
 \qquad
 \mathcal J^-_{ij}=\left.\frac{\partial f^-_j}{\partial u_i}\right|_{u=0}.
                                                                  \tag{KP.21}
\]
Consequently, at any integral fixed point
$\xi:B\to O_\xi$ with $O_\xi$ a finite coefficient DVR,
if $\mathcal J^-(\xi)$ has full row rank over its fraction
field, then
\[
 \operatorname{length}_{O_\xi}(N\otimes_{B,\xi}O_\xi)
       =v_{O_\xi}\!\left(\operatorname{Fitt}_0
                  (N\otimes_{B,\xi}O_\xi)\right)
       \leq v_{O_\xi}(\Delta_\xi)                            \tag{KP.22}
\]
for every nonzero maximal minor $\Delta_\xi$ of the evaluated
unsquared matrix.  The ideal in (KP.22) is independent of
the chosen presentation.
\end{lemma}
\begin{proof}
Write $J_0=J_0^+\oplus J_0^-$.  Commutativity and the two
projectors give
\[
              (J_0^2)^-=J_0^+J_0^- .                        \tag{KP.23}
\]
The ideal $Q$ is generated by $J_0^-$, its odd part is
$J_0^-$, and its even part consists of sums of products of
two odd elements, hence lies in $Q^2$.  Therefore $N$ is
entirely odd.  Also
$Q\cap J_0^2=J_0Q$: its odd part is (KP.23), and its even
part is contained in $Q^2\subset J_0Q$.  It follows that
\[
 Q/Q^2\otimes_B O=Q/J_0Q
       \simeq\operatorname{im}(Q\to J_0/J_0^2)
       =(J_0/J_0^2)^-,
\]
proving (KP.20).  This is an integral identity, not a
comparison with the reduced generic fixed locus.

In the displayed presentation, $Q$ is generated by the
classes of the $u_i$.  An even equation has no term linear
in the odd variables.  An odd equation, modulo $(u)^2$,
is exactly the column in (KP.21).  The even equations
specialized at $u=0$ define $B$.  These observations give
the right-exact presentation (KP.21), including all its
relations; multiples of the even equations vanish in $B$
and add none.  Base change to $O_\xi$ is right exact, and
Smith normal form proves (KP.22).  Presentation independence
is the defining invariance of the zeroth Fitting ideal.
\end{proof}

At an actual adjoint point, the coefficient module underlying
these odd normal cocycles has a particularly precise form.
Use left multiplicative cocycles
$r_\epsilon(g)=(1+\epsilon\gamma(g))r(g)$, exactly as in
\texttt{IsDeform}; thus their coefficient representation is
$\operatorname{ad}_{GG}r$, not its Tate dual.

\begin{lemma}[KT-ADJOINT-ODD-NORMAL-MODULE]
Let $\rho$ be the two-dimensional representation used in the
adjoint lift of KT-ADJOINT-ACTUAL-HOST-POINT, over a
characteristic-zero coefficient field.  For that literal
$\kappa$-fixed lift, the even and odd summands of
$\mathfrak{gl}_3$ as $G_{\mathbf Q}$-modules are
\[
 \begin{split}
   (\mathfrak{gl}_3)^+&\simeq\operatorname{Ad}^0\rho,\\
   (\mathfrak{gl}_3)^-&\simeq
       \eta\otimes\left(\mathbf1\oplus
                    \operatorname{Sym}^4\rho\otimes(\det\rho)^{-2}\right).
 \end{split}                                                     \tag{KP.24}
\]
Replacing $\rho$ by its dual leaves the normalized symmetric
fourth in this formula unchanged.
\end{lemma}
\begin{proof}
The differential of $\kappa$ is $X\mapsto-KX^tK$.
Its even space is $\mathfrak{so}(K)$ of dimension three,
and its odd space is the $K$-self-adjoint space of dimension
six.  Put $V=\operatorname{Ad}^0\rho$.  The invariant trace
form identifies $V$ with $V^*$, and the elementary
$\mathrm{GL}_2$ decomposition is
\[
 \bigwedge^2 V\simeq V,\qquad
 \operatorname{Sym}^2 V\simeq\mathbf1\oplus
                      \operatorname{Sym}^4\rho(\det\rho)^{-2}.
                                                                  \tag{KP.25}
\]
One way to verify the second identity is to write
$V=\operatorname{Sym}^2\rho\otimes(\det\rho)^{-1}$.
Multiplication of binary quadratics gives the surjection
$\operatorname{Sym}^2(\operatorname{Sym}^2\rho)
\to\operatorname{Sym}^4\rho$; its one-dimensional kernel is
$(\det\rho)^2$.  Over characteristic zero this sequence
splits.  The alternating identity follows from the
invariant volume and trace form of the determinant-one
three-dimensional representation.

The actual GG matrix is $A(g)D^{j(g)}$, with
$\eta(g)=(-1)^{j(g)}$ and $A(g)=\operatorname{Ad}^0\rho(g)$.
For $j(g)=1$ its action on the Lie algebra is
\[
   X\longmapsto A(g)D\theta_{\rm Lie}(X)D A(g)^{-1}
              =A(g)\kappa_{\rm Lie}(X)A(g)^{-1}.             \tag{KP.26}
\]
For $j(g)=0$ it is ordinary conjugation by $A(g)$.  Thus
the sign $\eta$ occurs on precisely the odd summand of
(KP.25), proving (KP.24).  Finally
$\rho^*\simeq\rho\otimes(\det\rho)^{-1}$ implies
\[
 \operatorname{Sym}^4(\rho^*)\otimes\det(\rho^*)^{-2}
       \simeq\operatorname{Sym}^4\rho\otimes(\det\rho)^{-2}.
\]
No Tate twist is introduced here.  A later use of local or
global duality would have to dualize and Tate-twist the
coefficient module separately.
\end{proof}

At a fixed two-sided point the carried characteristic-zero
rigidity statement \texttt{M4} leaves only the torus tangent.
That tangent is nonzero (some nonzero-weight function is
nonzero at the point) and is even because the point is fixed
and $\kappa_S$ commutes with the torus.  Hence every odd
normal tangent is zero.  Here the passage from the actual
ring to the finite graded model can be made explicitly.
A functional on $N\otimes_{B,\xi}\operatorname{Frac}O_\xi$
gives the odd derivation
$D(s)=\phi((s-\kappa_S(s))/2\bmod Q^2)$.
Since $N$ is a finite $B$-module and $\xi(B)\subset O_\xi$,
the values of $D$ are bounded in the fraction field.
After multiplication by some nonzero $c\in O$, the map
$f+\epsilon cD$ takes values in the complete coefficient
order
\[
 (O+\mathfrak m_{O_\xi})\oplus
                  \epsilon\mathfrak m_{O_\xi},\qquad\epsilon^2=0.
                                                                  \tag{KP.27}
\]
This order has the same residue field as $O$ and is
\texttt{IsCO}, by the same finite-order argument as in
KT-ADJOINT-ACTUAL-HOST-POINT.  If the restriction of $D$
to the model is zero, its value on every universal matrix
entry is zero.  Uniqueness in \texttt{S.rep}, applied to
(KP.27), makes the two homomorphisms $f+\epsilon cD$ and
$f$ equal; hence $D=0$.  The torus derivative is bounded
as well, by its functorial action over (KP.27).  Restrictions
therefore preserve linear independence between an odd
normal tangent and the even torus tangent, so \texttt{M4}
indeed excludes the former.  This avoids any assertion
about unrestricted algebraic differentials of a completed
power series ring.

It follows that (KP.21) has full generic row rank at such
a point, so (KP.22) is a finite local normal-length
invariant.  It does not make that invariant a unit:
characteristic-zero rigidity says only that the
corresponding integral normal module is torsion.

\paragraph{Precisely what remains for a branchwise replacement.}
The base inequality (KP.17) and the branch inequalities
(KP.22) are proved separately.  To replace an adjoint
trace-difference factor in the current global bound by
\emph{branchwise} normal Jacobian values, one still needs
an integral comparison bounding
$\operatorname{length}_O C^-_{-1}$ by those normal modules,
with appropriate residue degrees, ramification normalizations,
and a term controlling any vertical or embedded contribution.
The existing squared Rees construction (KP.2) does not give
that comparison.  A passage to the finite flat reduced
trace algebra does not give it either, by (KP.3)--(KP.4).
The exact missing assertion can also be stated as a kernel
condition: if a quotient $q:S\to S^{\rm flat,red}$ has trivial
involution, the image of the odd part of $\ker q$ in
$(J_0/J_0^2)_{-1}$ must be proved zero before that quotient
can discard the relevant Selmer conormal.  Point detection
alone proves no such containment in $J_0^2$.

The following conditional transfer makes the desired
numerical comparison concrete.  Its extra hypotheses are
not being asserted for the archived deformation model.

\begin{lemma}[KT-NORMAL-FLAT-TRANSFER, conditional criterion]
Let $B_0$ be a finite flat $O$-algebra whose generic fiber
is a product of fields $K_1,\ldots,K_s$.  Let $L$ be a
finite $B_0$-module equipped with a surjective $O$-linear map
$L\twoheadrightarrow C^-_{-1}$.  Suppose
\[
            B_0^m\xrightarrow{\mathcal J} B_0^n
                       \longrightarrow L\longrightarrow0
                                                                  \tag{KP.28}
\]
is a presentation, and $H$ is any $n$ by $n$ matrix obtained
by $B_0$-linear combinations of its relation columns such
that $\det H$ is nonzero in every $K_i$.  Let $O_i$ be the
integral closure of $O$ in $K_i$, let
$f_i=[k(O_i):k(O)]$, and let $H_i$ be the image of $H$ in
$M_n(O_i)$.  Then
\[
 \operatorname{length}_O\mathcal X
 \leq \operatorname{length}_O C^-_{-1}
 \leq \operatorname{length}_O L
 \leq v_O\!\left(\operatorname{Norm}_{B_0/O}(\det H)\right)
 =\sum_{i=1}^s f_i\,v_{O_i}(\det H_i).                      \tag{KP.29}
\]
In particular this replaces a trace comparison by normal
relation determinants whenever the stated integral
comparison is supplied.  Full row rank of $\mathcal J$
over every $K_i$ suffices for the existence of such an $H$.
\end{lemma}
\begin{proof}
The image of $H$ is contained in the image of
$\mathcal J$, so $\operatorname{coker}H$ surjects onto $L$.
Since $O$ is local, $B_0$ is finite free as an $O$-module.
After restricting scalars, $H$ is a square matrix over $O$
which is invertible over its fraction field.  Smith normal
form gives
\[
 \operatorname{length}_O\operatorname{coker}H
     =v_O(\det_O H)
     =v_O\!\left(\operatorname{Norm}_{B_0/O}(\det_{B_0}H)\right).
\]
The last equality can be checked after tensoring with the
fraction field and decomposing into the $K_i$; determinants
on a field extension restrict to field norms.  The usual
valuation of a field norm is $f_i$ times the normalized
valuation in $O_i$.  This proves the last equality of
(KP.29), and the quotient maps prove the two intermediate
inequalities.  Finite length now permits applying
KT-ODD-CONORMAL-GATE for the first one.

For the final assertion take $H=\mathcal J P$ with an
$m$ by $n$ matrix $P$ of indeterminates.  Full row rank in
$K_i$ makes $\det(\mathcal J_iP)$ a nonzero polynomial
for every $i$.  A finite union of their zero loci cannot
contain all matrices over the infinite field
$\operatorname{Frac}O$; equivalently expand coefficients
in bases of the finite extensions.  Choosing such a
matrix and clearing denominators supplies $P$ over $O$,
and hence the required $H$.
\end{proof}

For the actual application the full fixed-locus ring $B$
in (KP.19) retains torus directions and is not the finite
algebra required in this lemma.  One must descend the
relevant normal module to an appropriate finite
weight-zero algebra and control its integral torsion
before taking a flat reduced quotient.  A presentation
over that algebra, the surjection onto the base conormal,
and full generic rank on every retained component are
substantive hypotheses.  They are not consequences of
having a finite set of reduced automorphic isotypes.

Formula (KP.24) identifies the candidate arithmetic normal
deformation problem: the $\eta$-twisted normalized symmetric
fourth, together with the $\eta$ scalar line and the actual
ordinary/away-prime conditions.  An equality of its integral
Fitting ideal with an automorphic congruence ideal or a
special $L$-value would be an additional theorem.  None is
asserted here.  The present structural results are the
actual integral projector, its unsquared conormal presentation,
and the precise duality with the original strict Selmer
module.  No all-weight vanishing or removal of the adjoint
source family has been inferred from them.
